\subsection{User Flows}
\label{sec:user-flows}

The primary flow is: open Breakly $\rightarrow$ review the suggested
activity $\rightarrow$ start a break $\rightarrow$ follow the timer
$\rightarrow$ complete or skip $\rightarrow$ view the result. Supporting
flows cover onboarding, preference changes, reminder configuration, history,
offline use, and permission denial.

\subsection{Screen Inventory}
\label{sec:screen-inventory}

\begin{longtable}{|p{3cm}|p{0.31\textwidth}|p{0.39\textwidth}|}
\caption{Breakly screen inventory and state coverage}
\label{tab:screen-inventory}\\
\hline
\textbf{Screen} & \textbf{Purpose} & \textbf{Required states} \\
\hline
\endfirsthead
\multicolumn{3}{r}{\textit{Continued from previous page}}\\[-2pt]\hline
\textbf{Screen} & \textbf{Purpose} & \textbf{Required states} \\
\hline
\endhead
\multicolumn{3}{r}{\textit{Continue on next page}}\\
\endfoot
\endlastfoot
Home & Start the next break and see a recommendation & First use, ready,
offline, loading \\
\hline
Activity details & Explain the selected activity & Content, unavailable,
empty result \\
\hline
Active break & Show timer and controls & Running, paused, completed, skipped \\
\hline
History & Review previous sessions & Populated, empty, loading, error \\
\hline
Settings & Manage duration, reminders, and accessibility & Saved, validation
error, permission denied \\
\hline
\end{longtable}

\subsection{Visual and Accessibility Guidelines}
\label{sec:visual-accessibility}

Hi-fi mockups will use a consistent Material or iOS HIG-inspired component
system, spacing scale, typography hierarchy, and colour tokens. Every
relevant screen will specify loading, empty, offline, validation, and error
states. Accessibility review covers contrast, dynamic text size, screen-reader
labels, focus order, reduced motion, and touch targets of at least 44 dp.

\subsection{Prototype Evidence}
\label{sec:prototype}

The final submission will link a clickable prototype and include screenshots
of the primary and exception flows. Each screen in the prototype will carry
the same identifier used by the traceability matrix in
Chapter~\ref{chap:implementation-testing}.
